\section{Deciding Entailment in the Three Domains}
\label{sec:decision}

This section assembles the three bounds into decision procedures: a finite
search for holes, a procedure for arbitrary trees built on it
(\cref{fig:algorithm}), and transfers to regular and finite trees.

\begin{theorem}[Finite search for holes]\label{thm:hole-decision}
For an $N$-state reader, over any finite alphabet, that satisfies the
rejected-path property, put
\begin{equation}\label{eq:bounds}
 B=4N-2,\qquad d=2N(B+1)-1,\qquad H=\holeBound(N,d).
\end{equation}
A hole exists iff a hole of length at most $H$ exists, where $\holeBound$
is given by \eqref{eq:holebound}. Moreover,
$H\le2^{2^{1280(N+1)^5}}$.
\end{theorem}
\begin{proof}
By \cref{thm:brt}, choose a hole with tail at most $B$. If it is empty,
it already meets the bound. Otherwise, \cref{thm:depth} bounds its depth by
$d$ in the original $N$-state reader. Applying \cref{thm:search} gives a
hole of length at most $H$. For an empty alphabet only
$\eps$ exists, so no nonempty hierarchy is required. The reverse implication
is immediate. The closed-form estimate is \cref{lem:elementary-estimate}.
Each candidate can be checked by computing its run and
examining the finitely many pairs $s<e\le |w|$.
\end{proof}

\begin{figure}[t]
\small
\begin{minipage}{0.94\linewidth}
\textbf{Input:} arity $n$, variance $c$, a flat system $\phi$ on $k$ variables,
and query $x\le y$.\\
\textbf{Output:} whether $\phi\models_c x\le y$.
\begin{enumerate}[leftmargin=*,itemsep=2pt]
\item Form $\psi=\phi^\pm$, compute its closure, and test for a label clash.
If a clash exists, return \emph{yes}.
\item Instantiate the reader of \cref{sec:label-reader} on $\psi$ with the
roots $(x^+,y^+)$ and with the roots $(y^-,x^-)$.
Use $N=2\cdot4^{2k}$ and compute $B,d,H$ from \eqref{eq:bounds}.
\item Enumerate all words in $\Dir^*$ of length at most $H$, including
$\eps$. If either reader has a hole, return \emph{no}; otherwise return
\emph{yes}. (\Cref{sec:complexity} replaces this enumeration by the
summary search of \cref{prop:summary-search}.)
\end{enumerate}
For regular trees, use the same answer. For finite trees, first test finite
satisfiability; if it fails, return \emph{yes}, otherwise use the same answer.
\end{minipage}
\caption{The decision procedure for flat constraints and a variable query.}
\Description{A satisfiability test followed by two bounded searches for label-reader holes, with domain transfers for regular and finite trees.}
\label{fig:algorithm}
\end{figure}

\begin{theorem}[Correctness]\label{thm:correct}
The procedure of \cref{fig:algorithm} decides arbitrary-tree entailment.
\end{theorem}
\begin{proof}
By \cref{lem:selector}, \eqref{eq:g19}, and \cref{thm:sat}, step~1 returns yes exactly
when no solution exists. Otherwise, \eqref{eq:g13} and \cref{thm:holes} give
\[
 \phi\models_c x\le y\quad\Longleftrightarrow\quad
 \text{neither label reader has a hole}.
\]
Both readers satisfy the rejected-path property. By
\cref{thm:hole-decision}, the bounded search finds a hole whenever one exists.
Thus each answer is correct, and the finite enumeration terminates.
\end{proof}

\subsection{Regular and Finite Transfers}
Finite trees need one test beyond the label clash. A \emph{cycle clash} of
a covariant system $\chi$ is a nonempty word $\eta$ together with
variables $a,b$ such that
\begin{equation}\label{eq:g55}
 \Lo(\eta,a,a),\qquad D_\chi(a,b),\qquad \Up(\eta,b,b).
\end{equation}
In every solution $\sigma$ they give
$\emb_\bot(\eta,\sigma(a))\le\sigma(a)\le\sigma(b)\le\emb_\top(\eta,\sigma(b))$
by \eqref{eq:g23}. Iterating the two self-bounds forces constructor paths
along every power of $\eta$, so a cycle clash excludes finite solutions.

\begin{proposition}[Transfer]\label{prop:transfer}
For every arity and variance,
\begin{align}
 \phi\models_{\mathrm{reg},c}x\le y
 &\quad\Longleftrightarrow\quad\phi\models_c x\le y,\label{eq:g7}\\
 \phi\models_{\mathrm{fin},c}x\le y
 &\quad\Longleftrightarrow\quad
 \neg\Sat_{\mathrm{fin},c}(\phi)\ \lor\ \phi\models_c x\le y,\label{eq:g8}\\
 \Sat_{\mathrm{reg},c}(\phi)
 &\quad\Longleftrightarrow\quad\Sat_c(\phi).\label{eq:sat-reg}
\end{align}
Finite satisfiability is decidable.
\end{proposition}

\begin{proof}
The least shape is the unfolding of a finite graph of pairs of upper
and lower supporting-variable sets. Selection and decoding are realized
by graph products with polarity bits. Applying these operations to the
same four-spine extension as in \cref{thm:four-spine} gives a regular
countermodel whenever an arbitrary countermodel exists; the reverse
implication is immediate. The same operations applied to $\psi$ itself, without
spines, give a regular solution whenever a solution exists, which proves
\eqref{eq:sat-reg}.
The full graph construction and its equality
with the pruned least shape are proved in \cref{thm:shape-graph,thm:regular-transfer}.

For finite satisfiability, a covariant system must have neither a label
clash nor a cycle clash \eqref{eq:g55}. These conditions are sufficient by
\cref{thm:finite-sat}. Without a cycle clash the least shape has depth at
most $k_\chi^2+1$ for $k_\chi$ covariant variables; a longer constructor
path repeats a pair of supporting variables and yields a cycle clash.
A finite graph on such pairs decides the cycle test. Selection preserves
the finite domain, and decoding gives the finite signed correspondence
of \cref{lem:finite-signed}.

Finally, suppose the base has a finite solution and an arbitrary
countermodel. Its unsafe word has a clash-free four-spine extension.
Sources and sinks are separated. Assign old variables rank zero, new
spine variables their cut indices, and fillers rank $|w|+1$; ranks increase
along new constructor edges, so every cycle clash would reflect to the
base. The finite solution is used only to exclude cycle clashes in the
base. The extension preserves that exclusion, and its least shape can
therefore be reconstructed independently with bounded depth. Its selected fixed solution restricts and decodes to a
finite countermodel with the same unsafe word. This gives the finite
entailment equivalence, including the stated vacuity case; all rank and
depth details are proved in \cref{thm:finite-transfer}.
\end{proof}

The finite vacuity case is necessary: at positive arity the constraint
$f(x,\ldots,x)\le x\le f(x,\ldots,x)$ forces the infinite constructor
tree. It has no finite solution. For $n=0$, the reader search examines only
$\eps$, and the terminal child test still checks for a constructor bound.
Together with \cref{thm:correct}, the transfer proves \cref{thm:main}.

The effective graph-size and finite-depth bounds are collected in
\cref{sec:countermodel-bounds}.

The enumeration of \cref{fig:algorithm} establishes decidability.
\Cref{sec:complexity} implements its hole search differently: the summary
search (\cref{prop:summary-search}) stores one level of the block hierarchy
instead of a word and yields the space bound (\cref{thm:complexity}).
